\documentclass[../../../include/open-logic-section]{subfiles}

\begin{document}

\olfileid{sth}{ordinals}{ordtype}
\olsection{اعداد ترتیبی به‌مثابه نوع‌های ترتیب}

اکنون که با در اختیار داشتن جایگزینی در~$\ZFminus$ کار می‌کنیم، سرانجام
می‌توانیم نتیجه‌ای را اثبات کنیم که در پی آن بوده‌ایم:

\begin{thm}\ollabel{thmOrdinalRepresentation}
هر خوش‌ترتیب با یک عدد ترتیبی یکتا هم‌ریخت است.
\end{thm}

\begin{proof}
فرض کنید~$\tuple{A,<}$ یک خوش‌ترتیب باشد. بنا
بر~\olref[basic]{ordisoidentity}، این خوش‌ترتیب حداکثر با یک عدد ترتیبی
هم‌ریخت است. پس برای برهان خلف فرض کنید~$\tuple{A,<}$ با \emph{هیچ}
عدد ترتیبی‌ای هم‌ریخت نیست. ابتدا «$\tuple{A,<}$ را تا حد ممکن کوچک
می‌کنیم». به‌تفصیل: اگر قطعهٔ آغازین سره‌ای مانند~$\tuple{A_a,<_a}$
با هیچ عدد ترتیبی هم‌ریخت نباشد، عضو کمینه‌ای~$a\in A$ با این ویژگی
وجود دارد؛ در این صورت بگذارید~$B=A_a$ و
$\mathord{\lessdot}=\mathord{<_a}$. در غیر این صورت، بگذارید~$B=A$ و
$\mathord{\lessdot}=\mathord{<}$.

بنا بر تعریف، هر قطعهٔ آغازین سرهٔ~$B$ با عدد ترتیبی‌ای هم‌ریخت است که،
چنان‌که در بالا دیدیم، یکتاست. بنابراین با جایگزینی، مجموعهٔ زیر وجود دارد
و یک تابع است:
\[
	f = \Setabs{\tuple{\beta, b}}{b \in B\text{ و }\ordeq{\beta}{\tuple{B_b, \lessdot_b}}}
\]
برای تکمیل برهان خلف نشان می‌دهیم~$f$، برای یک عدد ترتیبی~$\alpha$،
هم‌ریختی‌ای~$\alpha\to B$ است.

آشکار است که~$\ran{f}=B$. همچنین بنا بر~\olref[iso]{lemordsegments}،
$f$ ترتیب را حفظ می‌کند؛ یعنی~$\gamma\in\beta$ اگر و تنها اگر
$f(\gamma)\lessdot f(\beta)$. برای نشان‌دادن اینکه~$\dom{f}$ عدد ترتیبی
است، بنا بر~\olref[basic]{corordtransitiveord} کافی است نشان دهیم~$\dom{f}$
متعدی است. پس~$\beta\in\dom{f}$ را ثابت بگیرید؛ یعنی برای یک~$b$،
$\ordeq{\beta}{\tuple{B_b, \lessdot_b}}$. اگر~$\gamma\in\beta$، آنگاه
بنا بر~\olref[iso]{wellordinitialsegment}، $\gamma\in\dom{f}$. با تعمیم،
$\beta\subseteq\dom{f}$.
\end{proof}

این نتیجه تعریف زیر را مجاز می‌سازد؛ تعریفی که از~\olref[vn]{sec} در پی
ارائهٔ آن بوده‌ایم:

\begin{defn}
اگر~$\tuple{A,<}$ یک خوش‌ترتیب باشد، نوع ترتیب آن، یعنی~$\ordtype{A,<}$،
عدد ترتیبی یکتای~$\alpha$ است به‌طوری‌که
$\ordeq{\tuple{A,<}}{\alpha}$.
\end{defn}

افزون بر این، این تعریف دو اصل خوشایند زیر را مجاز می‌سازد:

\begin{cor}\ollabel{ordtypesworklikeyouwant}
اگر~$\tuple{A,<}$ و~$\tuple{B,\lessdot}$ خوش‌ترتیب باشند:
\begin{align*}
	\ordtype{A, <} = \ordtype{B, \lessdot}&\text{ اگر و تنها اگر }\ordeq{\tuple{A, <}}{\tuple{B, \lessdot}}\\
	\ordtype{A, <} \in \ordtype{B, \lessdot}&\text{ اگر و تنها اگر }\ordeq{\tuple{A, <}}{\tuple{B_b, \lessdot_b}}\text{ برای یک }b \in B
\end{align*}
\end{cor}

\begin{proof}
تساوی بنا بر~\olref[basic]{ordisoidentity} برقرار است. برای اثبات ادعای
دوم، بگذارید~$\ordtype{A,<}=\alpha$ و~$\ordtype{B,\lessdot}=\beta$، و
بگذارید~$f\colon\beta\to B$ تابع زیربناییِ هم‌ریختی ترتیبی میان~$\beta$
و~$\tuple{B,\lessdot}$ باشد. آنگاه:
\begin{align*}
	\alpha \in \beta&\text{ اگر و تنها اگر }\funrestrictionto{f}{\alpha} \colon \alpha \to B_{f(\alpha)}\text{ یک هم‌ریختی باشد}\\
	&\text{ اگر و تنها اگر }\ordeq{\tuple{A, <}}{\tuple{B_{f(\alpha)}, \lessdot_{f(\alpha)}}}\\
	&\text{ اگر و تنها اگر }\ordeq{\tuple{A, <}}{\tuple{B_b, \lessdot_b}}\text{ برای یک $b \in B$}
\end{align*}
بنا بر~\olref[iso]{ordisounique}،
\olref[iso]{wellordinitialsegment} و
\olref[basic]{ordissetofsmallerord}.
\end{proof}

\end{document}
